\section{Reflection}

From the results of the test cases, the filter worked quite well for the given system parameters. For both tests, the maximum error did not exceed \textbf{5 cm}, while also maintaining yaw accuracy to within \textbf{0.5 rad}. Furthermore, while not shown in the figures, the velocity accuracy for both tests was within \textbf{0.35 m/s} maximum. Along with the incredible accuracy from both filters, the code is extremely efficient. On an Apple M2 Pro as a test bench, both test cases ran at \textbf{0.5 ms/predict \textit{and} update step}, over a period of 1500 iterations. As this framework was formulated specifically to run on less capable, miniature hardware such as an NVIDIA Jetson, this combination of accuracy and speed lends itself extremely well to a power-efficient localization application. Also, from all figures we notice that the $\pm 2\sigma$ error bounds for each parameter converges for all parameters - this is extremely important in the health of the filter. While the bounds for stationary parameters such as roll and pitch are a bit noisy, the bounds for yaw, and position in both test cases are quite smooth. This is a large indication that the orientation update and body kinematics are working in the way that we expect.

However, the current framework is not without its drawbacks. The main problem with the current dynamics formulation is that even with the two separate wheel ODEs from equations \eqref{eq:back_wheels} and \eqref{eq:front_wheels}, at slow RPMs the undriven wheels consistently switch directions, even though the friction coefficient $\mu$ is quite high, and the car is very unlikely to slip in the modeled conditions. Additionally, the oscillations shown in Figures \ref{fig:straight_rpy} and \ref{fig:sinusoidal_rpy} for the roll and pitch angles are likely due to the fact that the dynamics do not account for suspension, so the car is essentially a rigid body with no compliance, which is not entirely accurate. Furthermore, the current wheel dynamics are quite high frequency as well, with the current eigenvalues of the system exceeding $12$. This means that for traditional forward Euler, a timestep of \textit{167 \textmu s} is required to accurately resolve the dynamics, which is orders of magnitude too quick for accurate simulation. 

Future improvements on this project will focus on adding additional integration methods to ensure stability for a wider range of $\Delta t$ (i.e., backward Euler), while also further investigating the modeling of the wheels to find the root cause of negative RPMs while driving forward in nominal conditions.
